% =============================================================================
%  Capítulo 2 — Estado del Arte
%  Este capítulo es donde más texto tienes ya hecho. Trasvasa de tu PDF
%  actual aquí y aplica los arreglos puntuales del resumen de revisión.
% =============================================================================
\chapter{Estado del Arte}
\label{cap:estado-arte}

% -----------------------------------------------------------------------------
\section{Introducción al problema}
\label{sec:intro-problema}

\todo{Trasvasar de tu PDF actual la subsección "Introducción al problema",
"Objetivo" y "Estimación de densidad". Ya están bien redactadas.}

\subsection{Objetivo}
\label{sec:objetivo}
\todo{Texto existente.}

\subsection{Estimación de densidad}
\label{sec:estimacion-densidad}
\todo{Texto existente. Arreglar el cite \texttt{[2, Section 3.1]} si compila mal.}

% -----------------------------------------------------------------------------
\section{Modelos de difusión}
\label{sec:modelos-difusion}

\subsection{Proceso forward (inyección de ruido)}
\label{sec:forward}
\todo{Texto existente.}

\subsection{Score matching y función de pérdida}
\label{sec:score-matching}
\todo{Texto existente. \emph{Recordar:} añadir frase corta presentando la
U-Net como aproximador del score y remitir al apéndice
correspondiente.}

\subsection{Proceso backward (generación)}
\label{sec:backward}
\todo{Texto existente.}

\subsection{Variance Exploding (VE)}
\label{sec:ve}
\todo{Texto existente. Al final, añadir párrafo de \textbf{ventajas y
limitaciones} (lo pide el enunciado).}

\subsection{Variance Preserving (VP)}
\label{sec:vp}
\todo{Texto existente. Al final, párrafo de \textbf{ventajas y limitaciones}
y \textbf{tabla comparativa VE vs VP}.

\textbf{ARREGLAR}: el cosine schedule debe ir con cuadrado:
$f(t) = \cos^2\!\left(\frac{t/T+s}{1+s}\cdot\frac{\pi}{2}\right)$.}

\subsection{Ecuación de Fokker-Planck}
\label{sec:fp}
\todo{Texto existente.}

\subsection{Integración numérica de la SDE inversa: samplers}
\label{sec:samplers}

\todo{Subsección NUEVA. Tienes el archivo redactado en
\texttt{subseccion\_samplers.tex} (versión compacta). Pegar aquí.

Cubre los tres samplers (EM, PC, PF-ODE) que son requisito 2 del enunciado.}

% -----------------------------------------------------------------------------
\section{Generación controlable}
\label{sec:gen-controlable}

\todo{RENOMBRAR de "Modelos de procesos de difusión para generación
condicionada" a "Generación controlable" para que cuadre con el
enunciado (apartado 5).}

\subsection{El problema de la generación condicionada}
\todo{Texto existente.}

\subsection{Classifier guidance}
\todo{Texto existente.}

\subsection{Classifier-free guidance}
\todo{Texto existente.}

\subsection{Función de pérdida}
\todo{Texto existente.}

\subsection{Planteamiento de las SDEs}
\todo{Texto existente.}

\subsection{SDE backward condicionada}
\todo{Texto existente.}

\subsection{Imputación}
\label{sec:imputacion}
\todo{Subsección NUEVA. Tienes el archivo redactado en
\texttt{subseccion\_imputacion.tex}. Pegar aquí. Cubre el requisito 5b
del enunciado.}

% -----------------------------------------------------------------------------
\section{Métricas de calidad para las imágenes generadas}
\label{sec:metricas}

\subsection{Negative log-likelihood (NLL)}
\todo{Texto existente. \textbf{ARREGLAR}:
\begin{itemize}
  \item Eliminar la frase "(que para este problema pueden ser las mismas que
        para train)" — la NLL en train sobreestima.
  \item Nombrar explícitamente la \emph{dequantization} uniforme (Theis
        et al. 2016) en la definición $\v{y} = \v{x} + u$.
  \item La expresión del ELBO debe ser \emph{desigualdad} ($\geq$), no
        igualdad.
\end{itemize}}

\subsection{Bits per dimension (BPD)}
\todo{Texto existente.}

\subsection{Fréchet Inception Distance (FID)}
\todo{Texto existente.}

\subsection{Inception Score (IS)}
\todo{Texto existente. \textbf{ARREGLAR}: las referencias a las ecuaciones
(2.28) y (2.29) están rotas, no existen. Renumerar con \texttt{\textbackslash ref}.}

% -----------------------------------------------------------------------------
\section{Fenómeno del double descent}
\label{sec:double-descent}

\todo{O bien desarrollar como subsección (definición Belkin/Nakkiran 2019,
relación con la capacidad de la U-Net y el tamaño del dataset) o bien
\textbf{eliminar de aquí} y mover como observación empírica al
Capítulo de Conclusiones, citando los datos de §5.}
